\chapter*{Marco normativo}
\phantomsection
\addcontentsline{toc}{chapter}{Marco normativo}
\markboth{Marco normativo}{Marco normativo}

El proyecto utiliza conjuntos públicos de imágenes dermatológicas para investigación y prevé desarrollar un prototipo académico de apoyo al análisis. El prototipo aún no se ha implementado. Por ello, este apartado identifica normas y referencias relevantes para el tratamiento de información y una eventual etapa de uso, sin afirmar que exista aprobación ética institucional, autorización regulatoria o habilitación clínica.

\section*{Referentes internacionales}
\phantomsection
\addcontentsline{toc}{section}{Referentes internacionales}

El Reglamento General de Protección de Datos de la Unión Europea (GDPR) regula el tratamiento de datos personales en su ámbito de aplicación y clasifica los datos relativos a la salud y ciertos datos biométricos como categorías especiales. El tratamiento requiere una base jurídica conforme al Reglamento y, para datos de categorías especiales, una condición habilitante del artículo 9; el consentimiento explícito es una de las condiciones posibles, no la única. Las disposiciones sobre decisiones automatizadas se aplican a decisiones basadas únicamente en tratamiento automatizado que produzcan efectos jurídicos o afecten significativamente a una persona, con las condiciones y excepciones previstas en el Reglamento \parencite{ue2016gdpr}. La eventual aplicación del GDPR dependería de que se cumplan los criterios territoriales previstos en el propio Reglamento; aquí se incluye como referente comparativo.

La Declaración de Helsinki ofrece principios éticos para la investigación médica con participantes humanos. Se incluye como referente general; su aplicabilidad concreta a este proyecto, que trabaja con conjuntos públicos y no documenta recolección primaria, debe determinarse en la revisión institucional correspondiente \parencite{wma2024helsinki}.

El Reglamento de Inteligencia Artificial de la Unión Europea establece obligaciones para sistemas de alto riesgo en los casos definidos por el propio Reglamento, incluidos determinados sistemas que son productos sanitarios o componentes de seguridad sujetos a evaluación de conformidad. No todo sistema de IA usado para diagnóstico queda automáticamente clasificado como de alto riesgo. Se cita como referencia comparativa para una eventual comercialización o uso sujeto a jurisdicción europea, no como requisito vigente para este prototipo académico en Colombia \parencite{ue2024aiact}.

La Recomendación sobre la Ética de la Inteligencia Artificial de la UNESCO propone principios de derechos humanos, supervisión, transparencia, equidad y responsabilidad. Se considera una guía internacional y no una norma legal directamente vinculante para el proyecto \parencite{unesco2021eticaia}.

\section*{Referentes nacionales}
\phantomsection
\addcontentsline{toc}{section}{Referentes nacionales}

La Ley 1581 de 2012 establece el régimen general de protección de datos personales en Colombia y define como sensibles, entre otros, los datos relativos a la salud y los datos biométricos. Su tratamiento se rige por las bases, condiciones y excepciones previstas en la ley; por tanto, la disponibilidad pública de un conjunto no permite concluir, por sí sola, que cualquier reutilización o redistribución esté autorizada. Antes de una eventual aplicación con datos clínicos, deberán definirse las responsabilidades, la base jurídica y las salvaguardas correspondientes \parencite{colombia2012ley1581}.

La Ley 1419 de 2010 establece lineamientos para el desarrollo de la telesalud en Colombia. La Resolución 001644 de 2026 del Ministerio de Salud y Protección Social establece requisitos y condiciones para las actividades de telesalud y la prestación de servicios de telemedicina, y derogó la Resolución 2654 de 2019. Estas disposiciones se consideran relevantes si el prototipo llegara a integrarse en una prestación de servicios de salud; el presente trabajo no constituye un servicio de telemedicina ni declara cumplir requisitos de habilitación \parencite{colombia2010ley1419,minsalud2026resolucion1644}.

El Marco Ético para la Inteligencia Artificial en Colombia y el Documento CONPES 4144 ofrecen orientaciones de política y principios para el desarrollo y adopción responsable de IA. Se citan como referencias de contexto para discutir transparencia, responsabilidad y reducción de sesgos; no sustituyen la revisión normativa o ética aplicable al caso concreto \parencite{minciencias2021marcoetico,dnp2025conpes4144}.

\section*{Consideraciones éticas y alcance de uso}
\phantomsection
\addcontentsline{toc}{section}{Consideraciones éticas y alcance de uso}

La investigación documentada utiliza conjuntos públicos de imágenes y no reporta recolección primaria de datos de pacientes. Las licencias, condiciones de acceso y restricciones de redistribución deben verificarse para cada conjunto utilizado. La condición de repositorio público no demuestra por sí sola la aprobación de un comité de ética ni una exención institucional.

El prototipo funcional permanece como objetivo previsto y aún no se ha implementado ni desplegado. Antes de su desarrollo, deberán definirse la finalidad de uso, los datos procesados, la política de almacenamiento y eliminación, la información presentada al usuario, los mecanismos de supervisión humana y las responsabilidades asociadas. El sistema se plantea como apoyo experimental y no como sustituto del juicio profesional; no se reporta uso clínico ni certificación como dispositivo médico.

Los conjuntos disponibles no aportan una representación clínica suficiente para sostener conclusiones sobre todos los fototipos o sobre la población de Santander. La evaluación de tono realizada en el proyecto no valida el ITA como estimador de fototipo clínico. Estas limitaciones se deben reflejar en cualquier comunicación de resultados y en el diseño del prototipo.

La determinación del nivel de riesgo ético, los requisitos institucionales de la UNAB y la necesidad de revisión por una instancia ética quedan pendientes de consulta con el director de grado y la instancia correspondiente. Este documento no afirma que exista aprobación, exención o decisión institucional sobre esos aspectos.
